\section{Conclusions}\label{sec:conclusions}

We have presented a transient, differentiable inversion that recovers the distributed
surface mass balance of a glacier from two surface DEMs and a single surface-velocity
field, without ice-thickness measurements or per-glacier tuning. Surface velocity is
assimilated to fix a dynamically consistent state, and the mass-conservation equation
is then integrated forward with automatic differentiation to infer the
elevation-resolved SMB whose modelled geometry change best matches observation.
Applied through one glacier-independent pipeline to three Alpine glaciers of
contrasting size, the method reproduces in situ stake balances to
$0.32$--$0.65\,\mathrm{m\,w.e.\,yr^{-1}}$---on par with, and where GPR thickness is
available surpassing, the state-of-the-art
reconstruction of \citet{kneib2024}---and
a factor of two to five better than the
stationary-flux estimates it replaces.

The central advance is that mass conservation is enforced by construction: because the
flux divergence is recomputed from the evolving geometry at every step, the inferred
SMB needs neither the post-hoc smoothing nor the ad-hoc mass redistribution that
flux-divergence methods require, and stays dynamically consistent over the whole
observation window. Its accuracy is limited chiefly by the input velocity and DEMs and
by the one-dimensional elevation parametrisation, which trades sub-band spatial detail
for transferability. Because the method needs no calibration data beyond fields that
are approaching near-global coverage, it opens a realistic path towards continuous,
regional-scale SMB inversion; establishing its robustness across the smaller and more
varied glacier population that dominates regional water resources is the natural next
step.
